% Generated by build_report.py. Edit the shared content there, then regenerate.
\documentclass[10pt,a4paper]{article}
\usepackage{fontspec}
\setmainfont{texgyretermes}[Extension=.otf,UprightFont=*-regular,BoldFont=*-bold,ItalicFont=*-italic,BoldItalicFont=*-bolditalic]
\usepackage[margin=1.9cm,headheight=14pt]{geometry}
\usepackage{amsmath,amssymb,booktabs,tabularx,longtable,array,ragged2e,enumitem,microtype,xcolor,tikz,fancyhdr,graphicx}
\usetikzlibrary{arrows.meta,positioning}
\usepackage[hidelinks,unicode]{hyperref}
\definecolor{ink}{HTML}{183A45}\definecolor{teal}{HTML}{087F7A}\definecolor{pale}{HTML}{EAF4F2}\definecolor{gray}{HTML}{56656B}
\newcolumntype{P}[1]{>{\RaggedRight\arraybackslash}p{#1}}
\setlength{\parindent}{0pt}\setlength{\parskip}{5pt}\setlength{\emergencystretch}{2em}
\setlength{\tabcolsep}{4pt}\renewcommand{\arraystretch}{1.16}
\setlist{nosep,leftmargin=1.4em,topsep=4pt}
\pagestyle{fancy}\fancyhf{}\fancyhead[L]{\small\color{gray}BẢO VỆ RIÊNG TƯ QUỸ ĐẠO}\fancyhead[R]{\small\color{gray}03/10/2026}\fancyfoot[C]{\small\thepage}\renewcommand{\headrulewidth}{0.3pt}
\newcommand{\key}[1]{\par\smallskip\noindent\colorbox{pale}{\parbox{\dimexpr\linewidth-2\fboxsep}{\textbf{#1}}}\par\smallskip}
\hypersetup{pdftitle={Bảo vệ riêng tư quỹ đạo: dữ liệu, phương pháp và kết quả},pdfauthor={}}
\begin{document}
\begin{center}{\LARGE\bfseries\color{ink}Bảo vệ riêng tư quỹ đạo}\\[5pt]{\large Related works, kiến trúc và kết quả đối chứng}\\[4pt]{\small Bản trình bày rút gọn · 03/10/2026 · Số liệu giữ từ bản 30/09}\end{center}
\section{Related works ba năm gần đây và bảng metrics gốc}\label{sec:1}
Khảo sát đã cập nhật sang các nguồn trong \textbf{21/09/2023–21/09/2026} (mốc khóa khảo sát), tập trung vào bảo vệ vị trí/quỹ đạo và cách đánh giá. Bảng dưới rút gọn các metric đã kiểm tra trong hồ sơ nguồn; không xếp hạng bằng số liệu lấy từ các dataset khác nhau. CX = chưa xác minh đủ định nghĩa.

\begingroup\small
\begin{longtable}{@{}P{3.871cm}P{4.355cm}P{3.871cm}P{4.258cm}@{}}
\toprule
\textbf{Paper / hướng nghiên cứu} & \textbf{Privacy / chẩn đoán} & \textbf{Utility} & \textbf{Chi phí / phạm vi}\\\midrule
\endfirsthead
\toprule
\textbf{Paper / hướng nghiên cứu} & \textbf{Privacy / chẩn đoán} & \textbf{Utility} & \textbf{Chi phí / phạm vi}\\\midrule
\endhead
\bottomrule\endfoot
TransProtect 2024 [\hyperlink{ref-R1}{R1}] & EIE $\uparrow$: sai số suy luận & $\Delta$c $\downarrow$: chênh chi phí hành trình & Byte + latency đầu-cuối: CX\\
\addlinespace[3pt]
Semantic correlation 2026 [\hyperlink{ref-R2}{R2}] & ASR $\uparrow$: đạt ẩn danh; DER $\uparrow$: dummy/ngữ nghĩa & Chưa đo top-5 POI bằng DER & Thời gian sinh $\downarrow$\\
\addlinespace[3pt]
Fake queries 2026 [\hyperlink{ref-R3}{R3}] & Số đường khó phân biệt $\uparrow$; ASR $\uparrow$ & Chất lượng POI: CX & Thời gian sinh, RSS $\downarrow$\\
\addlinespace[3pt]
Road-aware PTPPM 2025 [\hyperlink{ref-R4}{R4}] & Sai số $\uparrow$; tấn công đúng $\downarrow$ & Dịch chuyển đầu ra $\downarrow$ & Chi phí LBS đầu-cuối: CX; preprint\\
\addlinespace[3pt]
CPCROK 2025 [\hyperlink{ref-R5}{R5}] & Tỷ lệ khớp quỹ đạo $\rho$ $\downarrow$ & Bối cảnh VANET & Số quỹ đạo giả $\psi$ $\downarrow$; proxy chi phí\\
\addlinespace[3pt]
DP-FETC 2025 [\hyperlink{ref-R6}{R6}] & MI $\downarrow$; bảo đảm $\varepsilon$-DP & JSD phân bố $\downarrow$ & Độ phức tạp; xuất bản offline\\
\addlinespace[3pt]
Improved PIR 2025 [\hyperlink{ref-R7}{R7}] & Phân tích an toàn; metric tấn công: CX & Truy hồi cụ thể: CX & Chi phí tính toán; mới đủ tóm tắt\\
\addlinespace[3pt]
SoK 2024 [\hyperlink{ref-R8}{R8}] & Tấn công tái dựng/liên kết & Phân bố, hình học, tác vụ & Khảo sát, không phải model\\
\addlinespace[3pt]
PRISM 2025 [\hyperlink{ref-R9}{R9}] & LPI; công thức CX & TDD, DC; định nghĩa CX & Thời gian xử lý/phản hồi\\
\addlinespace[3pt]
Options to Action 2026 [\hyperlink{ref-R10}{R10}] & Tỷ lệ dùng tính năng & Nhận thức và sử dụng & Nghiên cứu hành vi\\
\addlinespace[3pt]
Forsch et al. 2023 [\hyperlink{ref-R11}{R11}] & Trình bày S-TT / k-site-unidentifiability & Che điểm nhạy cảm & Chương tổng quan; kế thừa S-TT 2022\\
\addlinespace[3pt]
EV querying 2024 [\hyperlink{ref-R12}{R12}] & AGeoI + dummy & Dịch vụ truy vấn trạm sạc & Chưa xác minh benchmark endpoint\\
\end{longtable}
\endgroup
\textbf{Mốc nền giữ riêng:} DLS 2014 [\hyperlink{ref-B1}{B1}], RDG 2021 [\hyperlink{ref-B2}{B2}]; AnotherMe [\hyperlink{ref-B3}{B3}] có số tạp chí 2024 nhưng online 11/09/2023, ngoài cửa sổ. S-TT 2022 [\hyperlink{ref-B4}{B4}], tấn công EPZ 2022 [\hyperlink{ref-B5}{B5}] và Data Stream Obfuscator 07/2023 [\hyperlink{ref-B6}{B6}] bổ sung góc nhìn endpoint. Không gọi các nguồn này là paper mới trong ba năm.

\key{Kết quả khảo sát: metric gốc chấm các mục tiêu khác nhau; cần một bộ đo chung theo điều đối thủ suy ra, dịch vụ người dùng nhận và chi phí thực tế.}


\clearpage
\subsection{Vì sao chọn bộ metrics đề xuất?}\label{sub:1.1}
\begingroup\small
\begin{longtable}{@{}P{5.513cm}P{4.135cm}P{6.990cm}@{}}
\toprule
\textbf{Vấn đề từ related works} & \textbf{Metric đề xuất} & \textbf{Ý nghĩa trong bài toán của ta}\\\midrule
\endfirsthead
\toprule
\textbf{Vấn đề từ related works} & \textbf{Metric đề xuất} & \textbf{Ý nghĩa trong bài toán của ta}\\\midrule
\endhead
\bottomrule\endfoot
Entropy, DER, MI hoặc nhận dạng thật/ảo không cùng đáp án; ASR có nghĩa khác nhau giữa paper. & Hit50/100/200 $\downarrow$; MAE, median, p90 $\uparrow$ & Chấm tọa độ thật mà attacker suy ra, dù đầu ra là dummy hay vị trí thay thế. Hit100 là tỷ lệ đoán cách đáp án $\le$100 m.\\
\addlinespace[3pt]
MAE lớn vẫn có thể che một số lần suy đúng rất gần. & Đọc Hit cùng phân bố sai số & MAE cho độ lệch trung bình; Hit cho tỷ lệ lộ chính xác. S1: một điểm; S2: nơi dừng; S3: chuỗi điểm; S9/S10: endpoint.\\
\addlinespace[3pt]
JSD, $\Delta$c và DER chưa trả lời người dùng có nhận đúng POI đang khả dụng. & Recall@5 $\uparrow$; trả đủ $\uparrow$; $\Delta$D đường $\downarrow$ & So top-5 cục bộ với top-5 chuẩn dùng GPS thật và cùng trạng thái server. Ngưỡng Recall từng ca $\ge$90\%.\\
\addlinespace[3pt]
Số dummy/thời gian sinh bỏ sót phản hồi, query chèn và tải khi không có chuyến. & Byte yêu cầu + phản hồi $\downarrow$; latency p50/p95 $\downarrow$ & Đếm toàn bộ traffic trên mỗi sự kiện dịch vụ thật. Latency đầu-cuối chưa đo; chưa có điểm tổng Q thực nghiệm.\\
\end{longtable}
\endgroup
Recall@5 = số POI đúng có trong kết quả / số POI của tập chuẩn (tối đa 5). Yêu cầu có đáp án nhưng cơ chế không phục vụ nhận Recall=0; yêu cầu không có POI chuẩn để null. “Trả đủ” đo số lượng; $\Delta$D đo phần tăng khoảng cách đường. Hai chỉ số này bổ sung Recall khi phân tích utility chi tiết.

Để so được: giữ cùng mục tiêu, quyền quan sát, dịch vụ top-10/top-5 và quy tắc chọn attacker. Mỗi kiểu transcript có decoder phù hợp, được fit/chọn trên các nhóm riêng rồi khóa. Gộp seed/world trong record và nhóm, trung bình đều ca trong scenario rồi đều năm scenario; S10 chia hai ca A/B. Cùng dịch vụ vẫn phải báo riêng ngân sách truyền.

\key{Bộ đo đề xuất vận hành hóa ba câu hỏi: attacker suy đúng tới đâu; người dùng nhận được gì; phải trả bao nhiêu. Các metric là chỉ số đã có, được chọn và ghép cho nhiệm vụ này.}


\clearpage
\section{Kiến trúc mô hình BR-Dummy: Geo-I, mạng đường và các cơ chế}\label{sec:2}
\textbf{Nền tảng là Geo-I:} GPS thật tạo một neo riêng tư; từ lịch sử neo và ngữ cảnh công khai, mô hình sinh K vị trí/track giả có chuyển động hợp lệ và phủ POI hữu ích. Đây là kiến trúc cơ chế bảo vệ, với lớp boundary mở rộng cho điểm đầu/cuối.

\begin{center}
\includegraphics[width=\linewidth]{figures/architecture_model.pdf}
\end{center}
Đọc từ trái sang phải: input $\to$ layer bảo vệ Geo-I $\to$ layer ngữ cảnh $\to$ layer sinh dummy $\to$ output. Khung xanh bao toàn bộ mô hình của ta, gồm ba layer lõi và layer 4 boundary S9/S10 ở hàng dưới; input và bên nhận ở ngoài khung. Switching, pacing và slack tùy biến thể.

\begingroup\small
\begin{longtable}{@{}P{4.105cm}P{12.814cm}@{}}
\toprule
\textbf{Component} & \textbf{Cơ chế và vai trò đối với scenario}\\\midrule
\endfirsthead
\toprule
\textbf{Component} & \textbf{Cơ chế và vai trò đối với scenario}\\\midrule
\endhead
\bottomrule\endfoot
Geo-I / neo REM & Cơ chế mũ trên tập trạng thái đường công khai: xác suất giảm theo khoảng cách Euclid tới GPS. Tạo nền tảng bảo vệ vị trí S1; không cắt support theo bán kính bí mật quanh GPS.\\
\addlinespace[3pt]
Tái dùng neo + ngân sách & Test khoảng cách có nhiễu quyết định giữ/tạo neo; tính cả chi phí test và neo mới. Giảm nhiều mẫu nhiễu độc lập khi dừng (S2), kiểm soát tích lũy khi báo chuỗi (S3); pacing giãn đọc GPS, không đồng nghĩa che giờ chuyến.\\
\addlinespace[3pt]
Belief từ lịch sử neo & Ước lượng vùng vị trí/chuyển động chỉ từ neo đã bảo vệ; biến thể switching thêm chế độ dừng/đi. Hỗ trợ chọn dummy ở S2/S3, không đọc tốc độ thật hay đích tương lai.\\
\addlinespace[3pt]
Road-aware + POI-aware & Ràng buộc miền tới được trên làn có hướng/luật rẽ; chọn K điểm tối đa hóa độ phủ hợp POI, có exchange/slack tùy biến thể. Giữ chuỗi khả thi cho S3 và utility; chưa tự chứng minh chống attacker biết mạng đường.\\
\addlinespace[3pt]
S9: bỏ đầu trước lõi & BR-Boundary v1 không đưa GPS trong h giây đầu vào cơ chế neo, tránh dấu vết điểm đầu truyền vào trạng thái rồi lộ qua bản tin sau.\\
\addlinespace[3pt]
S10: buffer sau lõi & Giữ bản tin đã bảo vệ $\Delta$ giây; phát khi đủ tuổi; kết thúc thì hủy phần chưa phát. Không cần biết trước đích hoặc “60 giây cuối”. Có chi phí độ trễ và vẫn có rủi ro suy ngược từ biên còn thấy.\\
\end{longtable}
\endgroup
Sau neo, các khối belief/road/POI chỉ hậu xử lý dữ liệu đã bảo vệ và ngữ cảnh công khai, nên giữ cận lý tưởng của neo với composition đúng. K dummy không tự tạo k-anonymity. GPS thật còn dùng để lọc top-5 tại máy sau phản hồi; không đi vào selector. Lớp boundary đã có code và kiểm tra ghép lõi, nhưng chưa có benchmark đầu-cuối của toàn bộ Geo-I + boundary + dịch vụ.


\clearpage
\section{Benchmark: phân biệt lõi Geo-I và nhánh truy vấn công khai}\label{sec:3}
\key{Phạm vi bảng dưới: Đề xuất 30/67 là nhánh kế hoạch truy vấn công khai, không dùng neo Geo-I. Đây là kết quả của nhánh mở rộng dịch vụ; không gán Hit100=0 cho kiến trúc Geo-I/BR-Dummy hoặc BR-Boundary ở mục 2.}

Các đối chứng chạy cùng dịch vụ và được chấm trên cùng mục tiêu. \textbf{Năm adapter trực tuyến:} DLS, RDG, TransProtect*, Semantic*, Fake-query*. * TransProtect dùng Markov thay Transformer; Semantic dùng predictor thực nghiệm thay LSTM; Fake-query dùng giả định lịch chèn công khai. AnotherMe là tham chiếu offline, báo riêng phía dưới.

\subsection{Nhánh công khai: S1–S3 trên bốn nhóm tuyến}\label{sub:3.1}
54 record từ 30 chuyến trong bốn nhóm 901–904; attacker fit/chọn trên nhóm khác. Bảng là \textbf{Hit100 $\downarrow$}, trung bình đều A/B/C; utility ở p=0,8 (80\% khả dụng), trung bình đều năm scenario. Đây chưa phải kết quả của client theo lịch trên cohort 32 nhóm.

\begingroup\small
\begin{longtable}{@{}P{3.425cm}P{1.998cm}P{1.998cm}P{1.998cm}P{2.854cm}P{3.520cm}@{}}
\toprule
\textbf{Phương pháp} & \textbf{S1} & \textbf{S2} & \textbf{S3} & \textbf{Recall@5 $\uparrow$} & \textbf{Byte / sự kiện $\downarrow$}\\\midrule
\endfirsthead
\toprule
\textbf{Phương pháp} & \textbf{S1} & \textbf{S2} & \textbf{S3} & \textbf{Recall@5 $\uparrow$} & \textbf{Byte / sự kiện $\downarrow$}\\\midrule
\endhead
\bottomrule\endfoot
Vị trí thật & 100,00\% & 100,00\% & 100,00\% & 100,00\% & 242,8\\
\addlinespace[3pt]
DLS & 50,00\% & 41,67\% & 95,76\% & 100,00\% & 962,7\\
\addlinespace[3pt]
RDG & 58,33\% & 62,50\% & 95,83\% & 100,00\% & 950,7\\
\addlinespace[3pt]
TransProtect* & 66,67\% & 79,17\% & 74,59\% & 99,69\% & 242,7\\
\addlinespace[3pt]
Semantic* & 83,33\% & 100,00\% & 84,39\% & 100,00\% & 967,4\\
\addlinespace[3pt]
Fake-query* & 50,00\% & 33,33\% & 96,91\% & 100,00\% & 1.313,2\\
\addlinespace[3pt]
Đề xuất 30 & 0,00\% & 0,00\% & 0,00\% & 94,75\% & 2.389,7\\
\addlinespace[3pt]
Đề xuất 67 & 0,00\% & 0,00\% & 0,00\% & 100,00\% & 5.353,9\\
\end{longtable}
\endgroup
\textbf{Cách đọc:} vị trí thật và các adapter vẫn cho attacker suy đúng gần ở một phần ca; hai kế hoạch công khai đạt Hit100=0 trong bank đã thử. Điều này phù hợp với việc tọa độ công bố không bám GPS, nơi dừng hay tuyến thật. Không diễn giải 0\% thành an toàn trước mọi attacker.

\textbf{Utility và giá phải trả:} bản 30 giảm traffic so với 67 nhưng độ phủ POI thấp hơn. Bản 67 đạt Recall cao hơn nhờ thêm query; không coi privacy tốt hơn là thống trị cả ba trục. Byte là request + response JSON, giữ cả query phụ; chưa gồm HTTP/TLS. Bản 67 còn tám POI ngoài độ phủ tĩnh, nên Recall=100\% trên mẫu chưa là chứng nhận toàn bản đồ.


\clearpage
\subsection{Nhánh công khai: S9–S10 trên 32 nhóm tuyến, client theo lịch}\label{sub:3.2}
Cohort có 704 chuyến; utility dùng 435 record từ 246 chuyến, privacy endpoint dùng 150 record. Attacker khóa trước cohort này, bổ sung suy luận xuôi/ngược trên mạng đường và kết hợp các chuyến lặp. S9 gộp A/B/C, S10 gộp A/B; mỗi ô privacy là \textbf{Hit100 $\downarrow$ / MAE (m) $\uparrow$}.

\begingroup\small
\begin{longtable}{@{}P{3.177cm}P{3.658cm}P{3.658cm}P{2.695cm}P{2.888cm}@{}}
\toprule
\textbf{Phương pháp} & \textbf{S9: Hit / MAE} & \textbf{S10: Hit / MAE} & \textbf{Recall@5 $\uparrow$} & \textbf{Byte / sự kiện $\downarrow$}\\\midrule
\endfirsthead
\toprule
\textbf{Phương pháp} & \textbf{S9: Hit / MAE} & \textbf{S10: Hit / MAE} & \textbf{Recall@5 $\uparrow$} & \textbf{Byte / sự kiện $\downarrow$}\\\midrule
\endhead
\bottomrule\endfoot
Vị trí thật & 42,48\% / 165 & 17,39\% / 212 & 100,00\% & 249,9\\
\addlinespace[3pt]
DLS & 34,33\% / 462 & 11,16\% / 222 & 100,00\% & 952,2\\
\addlinespace[3pt]
RDG & 19,76\% / 793 & 10,75\% / 224 & 100,00\% & 944,4\\
\addlinespace[3pt]
TransProtect* & 33,61\% / 243 & 13,31\% / 259 & 99,70\% & 249,9\\
\addlinespace[3pt]
Semantic* & 40,81\% / 184 & 11,82\% / 292 & 100,00\% & 975,8\\
\addlinespace[3pt]
Fake-query* & 31,36\% / 522 & 26,11\% / 213 & 100,00\% & 1.295,6\\
\addlinespace[3pt]
30 + lịch & 0,00\% / 2.232 & 0,00\% / 1.916 & 95,00\% & 3.071,0\\
\addlinespace[3pt]
67 + lịch & 0,00\% / 2.311 & 0,00\% / 1.916 & 100,00\% & 6.878,6\\
\end{longtable}
\endgroup
\textbf{Privacy:} lịch công khai đạt Hit100=0 so với S9 \textbf{19,76–40,81\%} và S10 \textbf{10,75–26,11\%} của năm adapter. Raw control suy được endpoint ở cả S10.A (12,90\%) và S10.B (21,88\%), nên phép thử có sức phân biệt. CI 95\% của chênh lệch Hit100 nằm dưới 0 so với từng adapter ở cả S9/S10 (3.000 bootstrap theo nhóm; chưa hiệu chỉnh nhiều so sánh).

\textbf{Utility:} tại p=0,8, bản 30 đạt 95,00\%, 14/14 ca đạt $\ge$90\%; bản 67 đạt 100,00\%, 14/14 ca. Stress p=0,95: bản 30 đạt 93,01\%, 13/14 ca; bản 67 đạt 100\%, 14/14 ca ở các mức đã thử.

\textbf{Chi phí:} khoảng 3.071 / 6.879 byte mỗi sự kiện thật. So cùng kế hoạch chỉ refresh khi có hoạt động, lịch công khai tốn 4,48 lần byte: đây là giá của việc che giờ hoạt động, không phải tăng độ phủ POI.

\textbf{AnotherMe:} paper mô tả hệ thống online, nhưng adapter VTGA hiện có đọc cả chuyến. Trên 32 nhóm, privacy của tập con thành công có Hit100=0, Recall toàn phạm vi chỉ 16,85\%. Lỗi/thiếu đầu ra giữ utility bằng 0, privacy để thiếu; không xếp các số này như cùng mẫu số với năm adapter online.

\key{Kết luận cho nhánh công khai: thiết kế phối hợp kế hoạch công khai, lịch tải và xếp hạng cục bộ có lợi thế privacy trong năm scenario đã thử, đổi lại traffic cao hơn. Chưa chứng minh vượt các paper nguyên bản hoặc tốt hơn ở cùng ngân sách; dữ liệu vẫn cùng thành phố/bộ sinh SUMO.}


\clearpage
\section{Phụ lục: đọc chi tiết 14 sample A/B/C}\label{sec:4}
A/B/C là các điều kiện cùng nhiệm vụ, không phải thứ tự độ khó. Mỗi ca chọn một record thật từ bộ minh họa 12 nhóm/264 chuyến; số record không phải số chuyến độc lập. \textbf{S10 chỉ có A/B} trong tài liệu; B ánh xạ mã nguồn cũ S10.C để giữ truy nguyên, không có ca C độc lập trong phạm vi hiện tại.

Chấm màu = mẫu trước bảo vệ; nét đứt = đường thật để chấm; sao đỏ = đáp án; trục thời gian tách các mẫu chồng nhau. Attacker chỉ nhận transcript sau bảo vệ và phụ trợ được phép, không nhận các tọa độ/nhãn thật này. Map minh họa dùng nền SUMO cùng checksum OSM; thiếu mạng gốc để xác minh trùng toàn bộ hình học. © OpenStreetMap contributors.

\subsection{S1 — suy ra vị trí tại một lần gửi}\label{sub:4.1}
Attacker cần suy ra tọa độ tại một lần gửi. A/B thay ràng buộc mạng đường; C thêm ngữ cảnh POI hiếm.

\noindent\begin{minipage}[t]{0.48\linewidth}\vspace{0pt}\includegraphics[width=\linewidth]{figures/case_S1_A.pdf}\end{minipage}\hfill\begin{minipage}[t]{0.49\linewidth}\vspace{0pt}\RaggedRight \textbf{S1.A / 12 bản ghi.} Một bản tin tại cạnh có $\ge$2 hướng đi tiếp hợp lệ: kiểm tra khi mạng đường còn nhiều khả năng.\par\smallskip \textbf{Mẫu v2-r301-000.} u301\_\allowbreak{}00: 1 mẫu, FCD[333…333]\par\smallskip\textbf{Đọc sample.} Một lần lấy mẫu ở cạnh có 4 hướng đi tiếp. Sao đỏ trùng vị trí đầu vào; đối thủ phải chọn vị trí từ đầu ra đã bảo vệ, không được thấy sẵn sao này.\end{minipage}\par\vspace{3pt}
\noindent\begin{minipage}[t]{0.48\linewidth}\vspace{0pt}\includegraphics[width=\linewidth]{figures/case_S1_B.pdf}\end{minipage}\hfill\begin{minipage}[t]{0.49\linewidth}\vspace{0pt}\RaggedRight \textbf{S1.B / 12 bản ghi.} Một bản tin tại cạnh chỉ có 1 hướng đi tiếp: bản đồ tự loại bớt khả năng.\par\smallskip \textbf{Mẫu v2-r301-001.} u301\_\allowbreak{}00: 1 mẫu, FCD[241…241]\par\smallskip\textbf{Đọc sample.} Vẫn chỉ một lần lấy mẫu, nhưng cạnh chỉ có một hướng đi tiếp. Bản đồ có thể loại bớt khả năng so với A; cần đo lợi thế này bằng đối thủ thực tế.\end{minipage}\par\vspace{3pt}
\noindent\begin{minipage}[t]{0.48\linewidth}\vspace{0pt}\includegraphics[width=\linewidth]{figures/case_S1_C.pdf}\end{minipage}\hfill\begin{minipage}[t]{0.49\linewidth}\vspace{0pt}\RaggedRight \textbf{S1.C / 12 bản ghi.} Cách POI $\le$150 m; loại POI chiếm $\le$5\% danh mục công khai. Hiếm theo loại địa điểm, không phải theo tần suất ghé thăm.\par\smallskip \textbf{Mẫu v2-r301-029.} u301\_\allowbreak{}13: 1 mẫu, FCD[695…695]\par\smallskip\textbf{Đọc sample.} Điểm lấy mẫu cách POI pharmacy 12,03 m; loại này chiếm 3,11\% danh mục. Ngữ cảnh địa điểm hiếm là dấu hiệu bổ sung, không phải chuỗi di chuyển.\end{minipage}\par\vspace{3pt}

\clearpage
\subsection{S2 — suy ra nơi dừng qua nhiều lần gửi}\label{sub:4.2}
Attacker cần suy ra tọa độ nơi dừng. A/B thay thời lượng và nhịp quan sát; C kiểm tra rời đi rồi quay lại.

\noindent\begin{minipage}[t]{0.48\linewidth}\vspace{0pt}\includegraphics[width=\linewidth]{figures/case_S2_A.pdf}\end{minipage}\hfill\begin{minipage}[t]{0.49\linewidth}\vspace{0pt}\RaggedRight \textbf{S2.A / 12 bản ghi.} Dừng có chủ đích $\ge$20 giây; lấy truy vấn mỗi 5 giây.\par\smallskip \textbf{Mẫu v2-r301-002.} u301\_\allowbreak{}05: 6 mẫu, FCD[54…79]\par\smallskip\textbf{Đọc sample.} 6 mẫu cùng một tọa độ, cách nhau 5 giây; đoạn dừng thật dài 29 giây. Sáu vạch thời gian bị chồng thành một chấm trên map.\end{minipage}\par\vspace{3pt}
\noindent\begin{minipage}[t]{0.48\linewidth}\vspace{0pt}\includegraphics[width=\linewidth]{figures/case_S2_B.pdf}\end{minipage}\hfill\begin{minipage}[t]{0.49\linewidth}\vspace{0pt}\RaggedRight \textbf{S2.B / 12 bản ghi.} Dừng $\ge$120 giây; lấy truy vấn mỗi 20 giây.\par\smallskip \textbf{Mẫu v2-r301-003.} u301\_\allowbreak{}06: 9 mẫu, FCD[59…219]\par\smallskip\textbf{Đọc sample.} 9 mẫu cùng một tọa độ, cách nhau 20 giây; đoạn dừng dài 179 giây. Đối thủ có cửa sổ dài hơn để kết hợp bằng chứng, không chỉ một bản tin.\end{minipage}\par\vspace{3pt}
\noindent\begin{minipage}[t]{0.48\linewidth}\vspace{0pt}\includegraphics[width=\linewidth]{figures/case_S2_C.pdf}\end{minipage}\hfill\begin{minipage}[t]{0.49\linewidth}\vspace{0pt}\RaggedRight \textbf{S2.C / 12 bản ghi.} Hai lần dừng trên cùng làn, mỗi lần $\ge$20 giây; ở giữa thực sự di chuyển.\par\smallskip \textbf{Mẫu v2-r301-004.} u301\_\allowbreak{}07: 18 mẫu, FCD[214…1193]\par\smallskip\textbf{Đọc sample.} 18 mẫu chia hai đợt, mỗi đợt 9 mẫu. Hai lần dừng dài 44 giây tại cùng tọa độ; khoảng trống giữa hai cụm thời gian là lúc xe rời đi rồi quay lại.\end{minipage}\par\vspace{3pt}

\clearpage
\subsection{S3 — tái dựng đoạn đường đã đi}\label{sub:4.3}
Attacker cần tái dựng các vị trí của đoạn đã đi. A là chuỗi đều; B ít nhánh; C quan sát thưa.

\noindent\begin{minipage}[t]{0.48\linewidth}\vspace{0pt}\includegraphics[width=\linewidth]{figures/case_S3_A.pdf}\end{minipage}\hfill\begin{minipage}[t]{0.49\linewidth}\vspace{0pt}\RaggedRight \textbf{S3.A / 12 bản ghi.} Đoạn di chuyển $\ge$60 giây qua $\ge$3 cạnh; gửi mỗi 20 giây.\par\smallskip \textbf{Mẫu v2-r301-005.} u301\_\allowbreak{}00: 12 mẫu, FCD[361…581]\par\smallskip\textbf{Đọc sample.} 12 mẫu của u301\_\allowbreak{}00 cách nhau 20 giây. Chuỗi điểm cung cấp nhiều ràng buộc để nối tuyến; nét đứt là đường thật chỉ dùng chấm.\end{minipage}\par\vspace{3pt}
\noindent\begin{minipage}[t]{0.48\linewidth}\vspace{0pt}\includegraphics[width=\linewidth]{figures/case_S3_B.pdf}\end{minipage}\hfill\begin{minipage}[t]{0.49\linewidth}\vspace{0pt}\RaggedRight \textbf{S3.B / 11 bản ghi.} Ít nhất một nửa số cạnh được lấy mẫu chỉ có một hướng đi tiếp: hành lang ít lựa chọn.\par\smallskip \textbf{Mẫu v2-r301-007.} u301\_\allowbreak{}00: 9 mẫu, FCD[91…251]\par\smallskip\textbf{Đọc sample.} 9 mẫu; 55,6\% cạnh được lấy mẫu chỉ có một hướng đi tiếp. Đường ít nhánh có thể khiến chuỗi dễ suy hơn dù từng vị trí được làm mờ.\end{minipage}\par\vspace{3pt}
\noindent\begin{minipage}[t]{0.48\linewidth}\vspace{0pt}\includegraphics[width=\linewidth]{figures/case_S3_C.pdf}\end{minipage}\hfill\begin{minipage}[t]{0.49\linewidth}\vspace{0pt}\RaggedRight \textbf{S3.C / 12 bản ghi.} Lấy mẫu đoạn di chuyển thưa hơn, mỗi 60 giây.\par\smallskip \textbf{Mẫu v2-r301-006.} u301\_\allowbreak{}00: 4 mẫu, FCD[361…541]\par\smallskip\textbf{Đọc sample.} 4 mẫu của cùng đoạn trong A, cách nhau 60 giây. Khoảng trống lớn hơn nhưng đối thủ vẫn có thể dùng mạng đường để nối các khả năng.\end{minipage}\par\vspace{3pt}

\clearpage
\subsection{S9 — suy ra điểm xuất phát bị che}\label{sub:4.4}
Attacker cần suy điểm đầu bị giấu. A suy ngược một chuyến; B phân biệt hai nguồn nhập tuyến; C kết hợp chuyến lặp.

\noindent\begin{minipage}[t]{0.48\linewidth}\vspace{0pt}\includegraphics[width=\linewidth]{figures/case_S9_A.pdf}\end{minipage}\hfill\begin{minipage}[t]{0.49\linewidth}\vspace{0pt}\RaggedRight \textbf{S9.A / 12 bản ghi.} Cạnh xuất phát chỉ một lối ra. Che 60 giây đầu, lấy phần còn lại mỗi 20 giây.\par\smallskip \textbf{Mẫu v2-r301-025.} u301\_\allowbreak{}00: 27 mẫu, FCD[60…580]\par\smallskip\textbf{Đọc sample.} Sao là FCD[0]; cửa sổ bắt đầu ở FCD[60] sau khi che 60 giây. Chỉ một lối ra có thể giúp suy ngược nguồn xuất phát từ phần còn công bố.\end{minipage}\par\vspace{3pt}
\noindent\begin{minipage}[t]{0.48\linewidth}\vspace{0pt}\includegraphics[width=\linewidth]{figures/case_S9_B.pdf}\end{minipage}\hfill\begin{minipage}[t]{0.49\linewidth}\vspace{0pt}\RaggedRight \textbf{S9.B / 11 bản ghi.} Hai chuyến khác điểm đầu nhưng nhập một đoạn chung; chỉ cấp từ chỗ nhập tuyến.\par\smallskip \textbf{Mẫu v2-r301-026.} u301\_\allowbreak{}00: 28 mẫu, FCD[34…574]; u301\_\allowbreak{}04: 28 mẫu, FCD[19…559]\par\smallskip\textbf{Đọc sample.} Hai điểm đầu cách 281,76 m rồi nhập tuyến. Chỉ cấp dữ liệu sau nhập; hai sao là hai đáp án bị giấu mà đối thủ phải phân biệt.\end{minipage}\par\vspace{3pt}
\noindent\begin{minipage}[t]{0.48\linewidth}\vspace{0pt}\includegraphics[width=\linewidth]{figures/case_S9_C.pdf}\end{minipage}\hfill\begin{minipage}[t]{0.49\linewidth}\vspace{0pt}\RaggedRight \textbf{S9.C / 12 bản ghi.} Nhiều chuyến có cùng điểm xuất phát dự kiến; che 60 giây đầu từng chuyến.\par\smallskip \textbf{Mẫu v2-r301-027.} u301\_\allowbreak{}00: 27 mẫu, FCD[60…580]; u301\_\allowbreak{}09: 27 mẫu, FCD[60…580]\par\smallskip\textbf{Đọc sample.} u301\_\allowbreak{}00/u301\_\allowbreak{}09 lặp cùng điểm đầu dự kiến; mỗi phiên che 60 giây đầu. Nhiều cửa sổ có thể bổ sung bằng chứng về cùng nguồn xuất phát.\end{minipage}\par\vspace{3pt}

\clearpage
\subsection{S10 — suy ra điểm kết thúc bị che}\label{sub:4.5}
Attacker cần suy endpoint của chuyến đã hoàn tất. A dùng một chuyến; B kết hợp nhiều chuyến cùng đích. Đây khác dự đoán đích tương lai ở S6.

\noindent\begin{minipage}[t]{0.58\linewidth}\vspace{0pt}\includegraphics[width=\linewidth]{figures/case_S10_A.pdf}\end{minipage}\hfill\begin{minipage}[t]{0.39\linewidth}\vspace{0pt}\RaggedRight \textbf{S10.A / 11 bản ghi.} Cạnh kết thúc chỉ một lối vào. Che 60 giây cuối của chuyến đã hoàn tất.\par\smallskip \textbf{Mẫu v2-r302-025.} u302\_\allowbreak{}08: 38 mẫu, FCD[0…740]\par\smallskip\textbf{Đọc sample.} Sao là FCD cuối [817] của u302\_\allowbreak{}08; che 60 giây cuối. Đường chỉ một lối vào có thể giúp suy đoạn kết thúc từ phần trước còn thấy.\end{minipage}\par\vspace{3pt}
\noindent\begin{minipage}[t]{0.58\linewidth}\vspace{0pt}\includegraphics[width=\linewidth]{figures/case_S10_B.pdf}\end{minipage}\hfill\begin{minipage}[t]{0.39\linewidth}\vspace{0pt}\RaggedRight \textbf{S10.B / 12 bản ghi.} Nhiều chuyến cùng điểm kết thúc dự kiến; che 60 giây cuối từng chuyến.\par\smallskip \textbf{Mẫu v2-r301-028.} u301\_\allowbreak{}00: 27 mẫu, FCD[0…520]; u301\_\allowbreak{}09: 27 mẫu, FCD[0…520]\par\smallskip\textbf{Đọc sample.} u301\_\allowbreak{}00/u301\_\allowbreak{}09 lặp cùng đích dự kiến và cùng che 60 giây cuối. Đối thủ có thể kết hợp nhiều phiên để khoanh đích, chưa đủ gọi đó là nơi làm việc.\end{minipage}\par\vspace{3pt}
S9/S10 chấm tọa độ đầu/cuối; chưa gán nhãn nhà/nơi làm việc. Phóng to 29 ca của khung đầy đủ ở sample\_\allowbreak{}maps.html; tọa độ và record ở data\_\allowbreak{}samples.json, số đếm ở scenario\_\allowbreak{}inventory.csv.


\clearpage
\section{Nguồn và khả năng truy nguyên}\label{sec:5}
Bảng rút gọn đọc artifacts/benchmarks/active\_\allowbreak{}scope\_\allowbreak{}ac\_\allowbreak{}v2/readout.json đã xác minh, không chạy thêm model. method\_\allowbreak{}evidence.json giữ hash benchmark; preparation\_\allowbreak{}evidence.json giữ bảng guide. Hồ sơ metrics gốc nằm trong bản đầy đủ lưu ở archive/before\_\allowbreak{}concise\_\allowbreak{}2026-10-03/ và sources.json. Module concise\_\allowbreak{}presentation\_\allowbreak{}content.py tạo nội dung rút gọn; plot\_\allowbreak{}model\_\allowbreak{}architecture.py vẽ mô hình từ các component trong mã. scenario\_\allowbreak{}appendix.tex dùng chung cho hai tài liệu.

Ngày 03/10/2026 chỉ rút gọn nội dung và vẽ lại sơ đồ; mốc khảo sát và số liệu thực nghiệm được giữ riêng. Tra data\_\allowbreak{}samples.json để phân biệt nhãn trình bày với mã thực nghiệm, đặc biệt S10.B (mã nguồn S10.C cũ).

\par\hypertarget{ref-R1}{}\textbf{[R1]} Yadav et al. (2024). Protecting Vehicle Location Privacy with Contextually-Driven Synthetic Location Generation (TransProtect). SIGSPATIAL. \href{https://arxiv.org/html/2409.09495v1}{Nguồn gốc}. Toàn văn tác giả; §5.1.4, Eq. 13.
\par\hypertarget{ref-R2}{}\textbf{[R2]} Liu, Peng, Zhou (2026). A Dummy-Based Location Privacy Protection Scheme with Semantic Correlation of Moving Paths. JKSUCIS 38:478. \href{https://link.springer.com/article/10.1007/s44443-026-00899-w}{Nguồn gốc}. Toàn văn NXB; §6.3–6.5.
\par\hypertarget{ref-R3}{}\textbf{[R3]} Liu, Hu, Zhou (2026). Location Privacy Protection in Continuous LBSs: Enhancing Anonymity via Fake Queries. JKSUCIS 38:53. \href{https://link.springer.com/article/10.1007/s44443-025-00438-z}{Nguồn gốc}. Toàn văn NXB; §6.3–6.6.
\par\hypertarget{ref-R4}{}\textbf{[R4]} Road Network-Aware Personalized Trajectory Protection with Differential Privacy under Spatiotemporal Correlations (2025), arXiv:2511.21020v1. \href{https://arxiv.org/html/2511.21020v1}{Nguồn gốc}. Bản tiền công bố; §VI, Eq. 26–27.
\par\hypertarget{ref-R5}{}\textbf{[R5]} Wang et al. (2025). CPCROK: A Communication-Efficient and Privacy-Preserving Scheme for Low-Density Vehicular Ad Hoc Networks. Future Internet 17(4):165. \href{https://uhra.herts.ac.uk/id/eprint/25705/1/futureinternet-17-00165.pdf}{Nguồn gốc}. Toàn văn lưu tại trường tác giả; §5.4.
\par\hypertarget{ref-R6}{}\textbf{[R6]} Yue, Li, Liu, Li (2025). DP-FETC: A differentially private trajectory publishing method based on feature extraction and trajectory correlation. ERA 33(11):6631–6651. \href{https://www.aimspress.com/aimspress-data/era/2025/11/PDF/era-33-11-293.pdf}{Nguồn gốc}. Toàn văn NXB; §4.4, §5.2, Eq. 19–21.
\par\hypertarget{ref-R7}{}\textbf{[R7]} Gu et al. (2025). Research on Joint Protection of LBS Location and Query Privacy in Internet of Vehicles Based on an Improved PIR Algorithm. Concurrency and Computation: Practice and Experience. \href{https://onlinelibrary.wiley.com/doi/abs/10.1002/cpe.70447}{Nguồn gốc}. Chỉ xác minh tóm tắt NXB; chưa đủ định nghĩa các metric.
\par\hypertarget{ref-R8}{}\textbf{[R8]} Buchholz et al. (2024). SoK: Can Trajectory Generation Combine Privacy and Utility? PoPETs 2024(3):75–93. \href{https://petsymposium.org/popets/2024/popets-2024-0068.pdf}{Nguồn gốc}. Toàn văn; khảo sát và hướng dẫn đánh giá, không phải một cơ chế đối chứng.
\par\hypertarget{ref-R9}{}\textbf{[R9]} Farahnakiyan, Esmaeilyfard, Javidan (2025). A proactive privacy-preserving framework for mobile trajectory sharing (PRISM). JNCA 242:104271. \href{https://www.sciencedirect.com/science/article/pii/S1084804525001687}{Nguồn gốc}. Tóm tắt/đoạn công khai của NXB; chưa xác minh công thức LPI/TDD/DC.
\par\hypertarget{ref-R10}{}\textbf{[R10]} Ioannou, Aktypi, Athanasopoulos (2026). From Options to Action: Evaluating Adoption of Privacy Features in Fitness-Tracking Platforms. CHI 2026. \href{https://elathan.github.io/papers/chi26.pdf}{Nguồn gốc}. Toàn văn tác giả; bảng 3, §6.4. Đo sử dụng tính năng, không đo sức chống tấn công.
\par\hypertarget{ref-B1}{}\textbf{[B1]} Niu et al. (2014). Achieving k-Anonymity in Privacy-Aware Location-Based Services (DLS). INFOCOM. \href{https://doi.org/10.1109/INFOCOM.2014.6848002}{Nguồn gốc}. Đối chứng nền; định nghĩa entropy theo hồ sơ khảo sát đã có trong repo.
\par\hypertarget{ref-B2}{}\textbf{[B2]} Shaham et al. (2021). Privacy Preservation in Location-Based Services: A Novel Metric and Attack Model (RDG). IEEE TMC 20(10). \href{https://doi.org/10.1109/TMC.2020.2993599}{Nguồn gốc}. Đối chứng nền; entropy/chuyển tiếp/Viterbi theo hồ sơ đã có trong repo.
\par\hypertarget{ref-B3}{}\textbf{[B3]} Li et al. (2024 issue; online 11/09/2023). AnotherMe: A Location Privacy Protection System Based on Online Virtual Trajectory Generation. IEEE TDSC 21(4):2552–2567. \href{https://ieeexplore.ieee.org/document/10246991/}{Nguồn gốc}. Metadata/tóm tắt và mã nguồn tác giả; chưa kiểm chứng đầy đủ mẫu số/đơn vị phép đo.
\par\hypertarget{ref-B4}{}\textbf{[B4]} Brauer et al. (2022). My home is my secret: concealing sensitive locations by context-aware trajectory truncation. IJGIS 36(12):2496–2524. \href{https://doi.org/10.1080/13658816.2022.2081694}{Nguồn gốc}. Toàn văn NXB và chương tác giả 2023; S-TT, k-site-unidentifiability. Ngoại lệ nền trực tiếp cho S9/S10.
\par\hypertarget{ref-R11}{}\textbf{[R11]} Forsch et al. (2023). Efficient Mining of Volunteered Trajectory Datasets, §3.2. \href{https://link.springer.com/chapter/10.1007/978-3-031-35374-1_3}{Nguồn gốc}. Toàn văn tác giả, online 09/12/2023. Trình bày lại S-TT 2022, không phải một phương pháp mới độc lập.
\par\hypertarget{ref-B5}{}\textbf{[B5]} Dhondt et al. (2022). A Run a Day Won’t Keep the Hacker Away: Inference Attacks on Endpoint Privacy Zones in Fitness Tracking Social Networks. CCS. \href{https://people.cs.kuleuven.be/~stijn.volckaert/papers/2022_CCS_Fitness_Tracking.pdf}{Nguồn gốc}. Bản tác giả: tấn công EPZ và sáu countermeasures; không gắn nhãn một model bảo vệ đã thắng benchmark.
\par\hypertarget{ref-B6}{}\textbf{[B6]} Dehaene et al. (2023). Masking Location Streams in the Presence of Colluding Service Providers. EuroS\&P Workshops. \href{https://lirias.kuleuven.be/bitstream/20.500.12942/720391/2/Data_Stream_Obfuscator_camera_ready.pdf}{Nguồn gốc}. Bản tác giả: Data Stream Obfuscator, privacy zones/timeslots, mapping cache. Tháng 7/2023, ngoài cửa sổ ba năm.
\par\hypertarget{ref-R12}{}\textbf{[R12]} Atmaca et al. (2024). A Privacy-Preserving Querying Mechanism with High Utility for Electric Vehicles. IEEE OJVT. \href{https://wrap.warwick.ac.uk/id/eprint/183198/1/WRAP-privacy-preserving-querying-mechanism-high-utility-electric-vehicles-2024.pdf}{Nguồn gốc}. Bản tác giả: AGeoI + dummy cho truy vấn trạm sạc; không xác minh benchmark điểm đầu/cuối.
\end{document}
